%==============================================================================
%  PE3101P: Decision Theory and Social Choice
%  Take-Home Assignment 1 --- SUBMISSION TEMPLATE
%
%  HOW TO USE THIS FILE
%   1. Put your name and student number in the two lines just below
%      (between the braces).
%   2. Type each answer where it says "YOUR ANSWER".
%   3. Build with pdflatex (twice, so that cross-references settle), and
%      submit the PDF it produces.
%
%  You do not have to use this template. It is here to save you setting up a
%  document from scratch. You may delete anything you do not need.
%
%  If you have not used LaTeX before, you can open this file at overleaf.com
%  in a browser with nothing to install.
%==============================================================================
\documentclass[11pt]{article}

% ---- Fill these in ----------------------------------------------------------
\newcommand{\studentname}{YOUR NAME}
\newcommand{\studentnumber}{YOUR STUDENT NUMBER (A0\dots)}
% -----------------------------------------------------------------------------

\usepackage[utf8]{inputenc}
\usepackage[margin=1in]{geometry}
\usepackage{parskip}

\usepackage{amsmath}
\usepackage{amssymb}
\usepackage{nicefrac}
\usepackage{booktabs}
\usepackage{array}

\usepackage{enumitem}
\setlist{align=left}
\setlist[1]{labelindent=0em, itemsep=0em}
\setlist[itemize]{label=\textendash}
\setlist[enumerate,1]{label=\arabic*., itemsep=0.9em}
\setlist[enumerate,2]{label=(\alph*), itemsep=0.6em}
\setlist[enumerate,3]{label=(\roman*)}

\usepackage{amsthm}
\theoremstyle{definition}
\newtheorem*{definition*}{Definition}
\DeclareMathOperator{\Cr}{Cr}

\usepackage[hidelinks, colorlinks=true]{hyperref}
\hypersetup{linkcolor=black, citecolor=black, urlcolor=black}

% \begin{answer} ... \end{answer} marks where an answer goes. Write whatever
% suits the question: prose, displayed equations, a list of steps, a table.
\newenvironment{answer}
  {\par\medskip\noindent\textbf{Answer.}\par\nobreak\smallskip\ignorespaces}
  {\par\medskip}

\title{PE3101P: Decision Theory and Social Choice\\[0.3em]
       \large Take-Home Assignment 1\\[0.2em]
       \large Submission Template}
\author{\studentname \\ \studentnumber}
\date{Due Tuesday 8 September 2026, 23:59}

\begin{document}
\maketitle

\section*{Instructions}

Answer \textbf{all} questions. \textbf{The paper is marked out of 100.} The
marks available are shown against each question, in the right-hand margin.

You may use results proved in the lecture slides, in the lecture notes, or in
the primer \emph{Sets and Mathematical Proofs}, and you may use the probability
axioms freely, unless a question says otherwise. Where a question asks you to
prove something, set the proof out in steps and say what justifies each one.

This assignment is worth \textbf{20\%} of your final grade for the module.

%==============================================================================
\section*{Question 1 \hfill \normalsize [\,15 marks\,]}
%==============================================================================

Recall that, for any two acts $A$ and $B$ over the same set of states of nature:

\begin{itemize}
	\item[--] $A$ \textbf{strongly statewise dominates} $B$ if $A$ results in a
	better outcome than $B$ in \emph{every} state of nature.
	\item[--] $A$ \textbf{weakly statewise dominates} $B$ if $A$ results in an
	outcome at least as good as $B$'s in \emph{every} state of nature, and a
	better outcome in \emph{at least one} state of nature.
\end{itemize}

\begin{enumerate}[label=(\alph*)]

\item Show that if $A$ strongly statewise dominates $B$, and $B$ strongly
statewise dominates $C$, then $A$ strongly statewise dominates $C$.

\begin{answer}
% YOUR ANSWER to 1(a).
\end{answer}

\item Show that the same holds for weak statewise dominance: if $A$ weakly
statewise dominates $B$, and $B$ weakly statewise dominates $C$, then $A$ weakly
statewise dominates $C$.

\begin{answer}
% YOUR ANSWER to 1(b).
\end{answer}

\item Now suppose only one of the two links is strong. Show that if $A$ strongly
statewise dominates $B$ and $B$ weakly statewise dominates $C$, then $A$
\textbf{strongly} statewise dominates $C$; and that the same conclusion follows
if instead $A$ weakly statewise dominates $B$ and $B$ strongly statewise
dominates $C$.

\begin{answer}
% YOUR ANSWER to 1(c).
\end{answer}

\end{enumerate}

%==============================================================================
\section*{Question 2 \hfill \normalsize [\,10 marks\,]}
%==============================================================================

In each of the following, exactly one of the four options denotes the same event
as the expression at the top, whatever the events $A$, $B$ and $C$ happen to be.
Say which.

Throughout, $X^c$ is the complement of $X$, and $X \setminus Y$ is the set of
outcomes that are in $X$ but not in $Y$.

\begin{enumerate}[label=(\alph*)]

\item $(A \cup B)^c$
\begin{enumerate}[label=(\roman*), itemsep=0.1em, topsep=0.25em, parsep=0pt]
	\item $A^c \cup B^c$
	\item $A^c \cap B^c$
	\item $A^c \cup B$
	\item $(A \cap B)^c$
\end{enumerate}
\begin{answer}
% YOUR ANSWER to 2(a): write (i), (ii), (iii) or (iv).
\end{answer}

\item $A \cap (B \cup C)$
\begin{enumerate}[label=(\roman*), itemsep=0.1em, topsep=0.25em, parsep=0pt]
	\item $(A \cup B) \cap (A \cup C)$
	\item $(A \cap B) \cup C$
	\item $(A \cap B) \cup (A \cap C)$
	\item $A \cap B \cap C$
\end{enumerate}
\begin{answer}
% YOUR ANSWER to 2(b).
\end{answer}

\item $A \setminus (B \cup C)$
\begin{enumerate}[label=(\roman*), itemsep=0.1em, topsep=0.25em, parsep=0pt]
	\item $(A \setminus B) \cup (A \setminus C)$
	\item $(A \cup B) \setminus C$
	\item $A \setminus (B \cap C)$
	\item $(A \setminus B) \cap (A \setminus C)$
\end{enumerate}
\begin{answer}
% YOUR ANSWER to 2(c).
\end{answer}

\item $(A \cap B) \cup (A \cap B^c)$
\begin{enumerate}[label=(\roman*), itemsep=0.1em, topsep=0.25em, parsep=0pt]
	\item $B$
	\item $A$
	\item $A \cap B$
	\item $A \cup B$
\end{enumerate}
\begin{answer}
% YOUR ANSWER to 2(d).
\end{answer}

\item $(A \cup B) \cap (A \cup B^c)$
\begin{enumerate}[label=(\roman*), itemsep=0.1em, topsep=0.25em, parsep=0pt]
	\item $A \cup B$
	\item $\varnothing$
	\item $B$
	\item $A$
\end{enumerate}
\begin{answer}
% YOUR ANSWER to 2(e).
\end{answer}

\end{enumerate}

%==============================================================================
\section*{Question 3 \hfill \normalsize [\,15 marks\,]}
%==============================================================================

Recall the probability axioms. For any events $A$ and $B$:

\begin{itemize}
	\item[--] \textbf{Non-negativity:} $\Cr(A) \geq 0$.
	\item[--] \textbf{Normalisation:} $\Cr(\Omega) = 1$, where $\Omega$ is the
	sample space.
	\item[--] \textbf{Additivity:} if $A$ and $B$ are disjoint, then
	$\Cr(A \cup B) = \Cr(A) + \Cr(B)$.
\end{itemize}

In each of the following, an agent's credences are given. Say whether they can
all be held together without violating the axioms. If they cannot, say which
axiom or axioms are violated, and explain why. If one of two or more axioms must
be violated, but it is not settled which one, indicate this clearly.

\begin{enumerate}[label=(\alph*)]

\item A commuter is uncertain whether the 8:15 bus will arrive on time.
\[
	\Cr(\text{the bus is on time}) = 0.7,
	\qquad
	\Cr(\text{the bus is not on time}) = 0.4.
\]
\begin{answer}
% YOUR ANSWER to 3(a).
\end{answer}

\item A ball is to be drawn from a bag containing balls of several colours, some
of which are plastic.
\[
	\Cr(\text{the ball is red}) = 0.2,
	\qquad
	\Cr(\text{the ball is red and plastic}) = 0.35.
\]
\begin{answer}
% YOUR ANSWER to 3(b).
\end{answer}

\item A fair spinner is divided into eight equal sectors, numbered $1$ to $8$.
\[
	\Cr(\text{even}) = \nicefrac{1}{2},
	\quad
	\Cr(\text{lands on } 3) = \nicefrac{1}{8},
	\quad
	\Cr(\text{even, or lands on } 3) = \nicefrac{5}{8}.
\]
\begin{answer}
% YOUR ANSWER to 3(c).
\end{answer}

\item A football match is to be played. It can end in a home win, an away win,
or a draw.
\[
	\Cr(\text{a home win}) = 0.5,
	\quad
	\Cr(\text{a draw}) = 0.2,
	\quad
	\Cr(\text{a home win, or a draw}) = 0.8.
\]
\begin{answer}
% YOUR ANSWER to 3(d).
\end{answer}

\item A student is taking two modules this semester.
\[
	\Cr(\text{passes the first}) = 0.8,
	\quad
	\Cr(\text{passes the second}) = 0.7,
	\quad
	\Cr(\text{passes at least one}) = 0.9.
\]
\begin{answer}
% YOUR ANSWER to 3(e).
\end{answer}

\end{enumerate}

%==============================================================================
\section*{Question 4 \hfill \normalsize [\,15 marks\,]}
%==============================================================================

Let $\Cr$ be a credence function satisfying the probability axioms. Suppose the
sample space $\Omega$ can be partitioned into $B_1, B_2, \dots, B_n$ --- that is,
the $B_i$ are mutually exclusive and jointly exhaustive. Let $A$ be any event.

\begin{enumerate}[label=(\alph*)]

\item Show that
\[
	\Cr(A) = \Cr(A \cap B_1) + \Cr(A \cap B_2) + \dots + \Cr(A \cap B_n).
\]

\emph{You may assume the following, which you are not required to prove: if $B$
and $C$ are disjoint, then $A \cap B$ and $A \cap C$ are disjoint. You may also
assume that $A = (A \cap B_1) \cup (A \cap B_2) \cup \dots \cup (A \cap B_n)$.}

\begin{answer}
% YOUR ANSWER to 4(a).
\end{answer}

\item Hence prove the \textbf{law of total probability}: where $\Cr(B_i) > 0$
for every $i$,
\[
	\Cr(A) = \Cr(B_1)\Cr(A \mid B_1) + \Cr(B_2)\Cr(A \mid B_2)
	         + \dots + \Cr(B_n)\Cr(A \mid B_n).
\]

\begin{answer}
% YOUR ANSWER to 4(b).
\end{answer}

\end{enumerate}

%==============================================================================
\section*{Question 5 \hfill \normalsize [\,15 marks\,]}
%==============================================================================

Recall the definition of conditional probability: for any events $X$ and $Y$,
\[
	\Cr(X \mid Y) = \frac{\Cr(X \cap Y)}{\Cr(Y)}.
\]
Assume throughout this question that every event mentioned has positive credence.

\begin{enumerate}[label=(\alph*)]

\item Using only this definition, prove \textbf{Bayes' theorem}:
\[
	\Cr(A \mid B) = \frac{\Cr(B \mid A) \cdot \Cr(A)}{\Cr(B)}.
\]

\begin{answer}
% YOUR ANSWER to 5(a).
\end{answer}

\item Using your answer to part (a), show that
\[
	\underbrace{\frac{\Cr(H_1 \mid E)}{\Cr(H_2 \mid E)}}_{\text{posterior odds}}
	\;=\;
	\underbrace{\frac{\Cr(E \mid H_1)}{\Cr(E \mid H_2)}}_{\text{Bayes factor}}
	\;\times\;
	\underbrace{\frac{\Cr(H_1)}{\Cr(H_2)}}_{\text{prior odds}}.
\]

\begin{answer}
% YOUR ANSWER to 5(b).
\end{answer}

\end{enumerate}

%==============================================================================
\section*{Question 6 \hfill \normalsize [\,15 marks\,]}
%==============================================================================

A bag contains three six-sided dice, which look and feel identical. One of them
is loaded,
and lands six with probability $\nicefrac{1}{2}$; the other two are fair. You
take one die from the bag at random, roll it once, and it lands six.

Let $H_1$ be the hypothesis that you took the loaded die, $H_2$ the hypothesis
that you took a fair one, and let $E$ be the evidence that the roll landed six.

\begin{enumerate}[label=(\alph*)]

\item Write down the prior odds of $H_1$ over $H_2$, and the Bayes factor.

\begin{answer}
% YOUR ANSWER to 6(a).
\end{answer}

\item Using the odds form of Bayes' theorem, as stated in question 5(b),
calculate the posterior odds, and hence your credence in $H_1$ and in $H_2$ after
seeing the roll.

\begin{answer}
% YOUR ANSWER to 6(b).
\end{answer}

\item Suppose instead that the roll had \emph{not} landed six, and everything
else had gone as before. What would your credence in $H_1$ have been? Which of
the prior odds, the Bayes factor and the posterior odds changed, and which did
not?

\begin{answer}
% YOUR ANSWER to 6(c).
\end{answer}

\end{enumerate}

%==============================================================================
\section*{Question 7 \hfill \normalsize [\,15 marks\,]}
%==============================================================================

In a decision problem with finitely many acts and finitely many possible outcomes,
a \textbf{dependency hypothesis} lists, for each act available, which outcome
would result were that act performed.

It will help to have a compact notation. Where there are two acts, write
$\langle o, o' \rangle$ for the dependency hypothesis on which the first act would
produce outcome $o$ and the second would produce outcome $o'$; and similarly, with
a third entry, where there are three acts. Each outcome below is given a letter to
use in place of its name.

\emph{For example}, take the acts \emph{revise} and \emph{go out}, with outcomes
\emph{pass} (p) and \emph{fail} (f). One of the dependency hypotheses is
$\langle p, f \rangle$: revising would result in passing, and going out would
result in failing.

\begin{enumerate}[label=(\alph*)]

\item In each of the following, list \textbf{every} possible dependency
hypothesis.

\begin{enumerate}[label=(\roman*), itemsep=0.1em, topsep=0.25em, parsep=0pt]

\item Acts: \emph{take the umbrella}, \emph{leave the umbrella}.
Outcomes: \emph{dry} (d), \emph{soaked} (s).

\begin{answer}
% YOUR ANSWER to 7(a)(i).
\end{answer}

\item Acts: \emph{water the plant}, \emph{leave the plant}.
Outcomes: \emph{flowers} (f), \emph{survives} (s), \emph{dies} (d).

\begin{answer}
% YOUR ANSWER to 7(a)(ii).
\end{answer}

\item Acts: \emph{brush the cat}, \emph{stroke the cat}, \emph{leave the cat
alone}. Outcomes: \emph{purrs} (p), \emph{quiet} (q).

\begin{answer}
% YOUR ANSWER to 7(a)(iii).
\end{answer}

\end{enumerate}

\item In general, how many dependency hypotheses are there when a decision problem
has $m$ acts and $n$ possible outcomes? Explain your answer.

\begin{answer}
% YOUR ANSWER to 7(b).
\end{answer}

\end{enumerate}

\end{document}
